%===========================================
% 第10章 考察
% 推定結果の解釈、四要素の相対的重要性、政策的合意を述べる。Zlinの構造的特徴と新電力市場の制度的脆弱性の関係を検討する。
%===========================================

\chapter{考察}
\label{chap:discussion}


\section{政策的含意}

本研究は，新電力企業の倒産が単発の価格ショックではなく，
ビジネスモデルの構造的脆弱性によって決定されることを明らかにした。
本節では，得られた知見を踏まえ，新電力制度，とりわけ新規参入政策に対する
政策的含意を整理する。

\subsection{ヘッジ市場の整備と先渡し契約の普及}

本研究の分析により，外部ショック感応度 $\sigma$ が倒産確率を大きく押し上げる
主要因であることが示された。これは，スポット市場への依存度が高い企業ほど，
短期的な価格変動を吸収できず，逆ザヤによって資金繰りが急速に悪化することを意味する。

したがって，政策的には以下が重要となる。

\begin{itemize}
    \item 先渡し契約（ヘッジ）の最低比率の設定
    \item 長期契約市場の流動性向上
    \item 小規模事業者向けのヘッジ支援スキームの創設
\end{itemize}

欧州の電力市場では，ヘッジ比率の義務化や長期契約市場の整備が進んでおり，
これらは新電力の $\sigma$ を制度的に低減させる政策として有効である。

\subsection{支援能力 $S$ を高める制度設計}

本研究では，支援能力 $S$ が倒産確率を大幅に低下させることが示された。
親会社・自治体・金融機関などの外生的支援主体の存在は，
短期的な逆ザヤや制度コストの急増を吸収する「緩衝材」として機能する。

政策的には，以下の制度設計が考えられる。

\begin{itemize}
    \item 親会社保証の制度化（一定規模以下の新電力に対する義務化）
    \item 自治体新電力に対する財政的・制度的支援の明確化
    \item 金融機関による信用保証枠の創設
\end{itemize}

特に小規模新電力にとっては，$S$ の有無が倒産確率を決定づけるため，
支援能力の制度的確保は新規参入政策の中核となる。

\subsection{倒産閾値 $D$ を引き下げる制度改革}

倒産閾値 $D$ は，インバランス負担や調整市場コストなど，
外部ショック時に企業が最低限支払わなければならないキャッシュアウトを表す。
本研究の結果は，$D$ の上昇が倒産確率を急激に押し上げることを示している。

したがって，以下の制度改革が政策的含意として導かれる。

\begin{itemize}
    \item インバランス制度の見直し（小規模事業者への負担軽減）
    \item 需給調整市場の設計改善（過度なコスト転嫁の抑制）
    \item JEPX 担保制度の柔軟化（小規模事業者向けの担保要件緩和）
\end{itemize}

これらは，$D$ を制度的に引き下げることで，
新電力の倒産リスクを構造的に低減させる政策である。

\subsection{新規参入政策の再設計}

本研究の結果は，新規参入政策に対しても重要な示唆を与える。
従来の新規参入政策は，参入障壁の低減を中心として設計されてきたが，
構造的脆弱性を抱えたままの参入は，市場の不安定化を招く可能性がある。

したがって，新規参入に際しては以下の観点が必要となる。

\begin{itemize}
    \item 参入時における最低限のヘッジ比率の確認
    \item 財務的耐性（$\mu$）および支援能力（$S$）の事前評価
    \item 需給予測能力や調達戦略の審査
\end{itemize}

これらは参入障壁を高めるものではなく，
市場の持続可能性を確保するための「参入要件の適正化」である。

\subsection{政策的含意のまとめ}

本研究の構造モデル（$\mu, \sigma, S, D$）は，
新電力の倒産リスクが財務指標ではなく，
市場構造・制度構造・支援構造によって決定されることを示した。

したがって，政策的には，

\begin{itemize}
    \item $\sigma$ を下げる（ヘッジ市場の整備）
    \item $S$ を上げる（支援能力の制度化）
    \item $D$ を下げる（制度コストの見直し）
\end{itemize}

という三方向の改革が必要である。

これらの政策的含意は，本研究の理論モデルと実証結果の双方に基づくものであり，
新電力市場の持続可能性を高めるための実務的かつ現実的な提案である。
